\lecture{6.4}{Basic Calls to \texttt{SuperLearner}}
\label{lec:6-4}

\subsection{Main Arguments: \texttt{Y}, \texttt{X}, \texttt{SL.library}, \texttt{family}, \texttt{method}, \texttt{cvControl}}

The function \texttt{SuperLearner()} automates the cross-validation and ensemble
construction described in the previous lectures. Its most important arguments are:
\begin{itemize}
\item \texttt{Y}: the outcome to be predicted;
\item \texttt{X}: a data frame or matrix of predictors;
\item \texttt{SL.library}: candidate algorithms;
\item \texttt{family}: the outcome family, commonly \texttt{gaussian()} or
\texttt{binomial()};
\item \texttt{method}: the metalearner and loss used to choose weights;
\item \texttt{cvControl}: options controlling the internal cross-validation.
\end{itemize}

For outcome regression, \texttt{Y} is the observed outcome and \texttt{X} contains
$(A,W)$. For propensity score estimation, \texttt{Y} is the treatment $A$ and
\texttt{X} contains $W$.

\begin{implementationbox}
\begin{lstlisting}[language=R]
set.seed(123)

sl_fit1 <- SuperLearner(
  Y = Y,
  X = data.frame(W1 = W1, W2 = W2, A = A),
  SL.library = c("SL.glm", "SL.mean"),
  family = gaussian(),
  method = "method.CC_LS",
  cvControl = list(V = 2),
  verbose = FALSE
)
\end{lstlisting}
\end{implementationbox}

\subsection{Continuous vs Binary Outcomes}

For a continuous outcome regression, use \texttt{family = gaussian()} and a loss such
as squared-error loss. A common metalearner is \texttt{method.CC\_LS}, a convex
combination chosen to minimize least-squares loss.

For a binary outcome or binary treatment mechanism, use \texttt{family = binomial()}.
For predicted probabilities, a natural loss is negative log-likelihood. A useful
metalearner is \texttt{method.CC\_nloglik}.

\begin{implementationbox}
\begin{lstlisting}[language=R]
# Outcome regression for continuous Y
fit_Q <- SuperLearner(Y = Y,
                      X = data.frame(A = A, W1 = W1, W2 = W2),
                      SL.library = c("SL.glm", "SL.mean"),
                      family = gaussian(),
                      method = "method.CC_LS")

# Propensity score for binary A
fit_g <- SuperLearner(Y = A,
                      X = data.frame(W1 = W1, W2 = W2),
                      SL.library = c("SL.glm", "SL.mean"),
                      family = binomial(),
                      method = "method.CC_nloglik")
\end{lstlisting}
\end{implementationbox}

\subsection{Interpreting Printed Output}

Printing a \texttt{SuperLearner} object shows the call, the cross-validated risk of
each candidate, and the final ensemble coefficients. A typical table has columns:
\[
\text{Risk},\qquad \text{Coef}.
\]
The risk column describes candidate performance under the chosen cross-validated
loss. The coefficient column gives the ensemble weight assigned to each candidate.

\subsection{Risk}

The reported risk is not the full-data training error. It is an internal
cross-validated risk estimate for each candidate algorithm. For Gaussian outcomes
with \texttt{method.CC\_LS}, this risk is mean squared prediction error. For binomial
outcomes with \texttt{method.CC\_nloglik}, it is negative log-likelihood risk.

Risk values are comparable only when they use the same outcome and same loss. A
mean squared-error risk for $Y$ should not be numerically compared with a
negative-log-likelihood risk for $A$.

\subsection{Coefficients}

The coefficients are the estimated ensemble weights. If a candidate has coefficient
zero, it does not contribute to the final Super Learner prediction. This does not mean
the candidate is scientifically impossible or universally poor; it means that, given the
other candidates and the selected loss, the metalearner did not need that candidate to
reduce cross-validated risk.

\begin{keyidea}
The risk table scores individual algorithms. The coefficients describe the ensemble.
An algorithm can have mediocre individual risk but still receive weight if its
predictions complement the predictions of other learners.
\end{keyidea}

\subsection{Super Learner Predictions}

The object returned by \texttt{SuperLearner()} contains fitted predictions for the
observed rows:

\begin{implementationbox}
\begin{lstlisting}[language=R]
head(sl_fit1$SL.predict)
\end{lstlisting}
\end{implementationbox}

Predictions for new observations are obtained with \texttt{predict()}:

\begin{implementationbox}
\begin{lstlisting}[language=R]
new_obs <- data.frame(A = 1, W1 = 0, W2 = 0)

new_predictions <- predict(sl_fit1, newdata = new_obs)
new_predictions$pred
\end{lstlisting}
\end{implementationbox}

The returned \texttt{pred} component is the ensemble prediction.

\subsection{Library Predictions}

The same prediction call also returns predictions from each library member:

\begin{implementationbox}
\begin{lstlisting}[language=R]
new_predictions$library.predict
\end{lstlisting}
\end{implementationbox}

Inside the fitted object, useful components include:
\begin{itemize}
\item \texttt{SL.predict}: Super Learner predictions in the original data;
\item \texttt{library.predict}: full-data predictions from each candidate;
\item \texttt{coef}: ensemble weights;
\item \texttt{cvRisk}: internal cross-validated risk estimates;
\item \texttt{fitLibrary}: fitted candidate algorithms;
\item \texttt{validRows}: validation fold membership.
\end{itemize}

\begin{estimationbox}
For causal estimation, the fitted Super Learner is still a nuisance-function estimator.
To use it in g-computation, predict $\Qbar_n(1,W_i)$ and $\Qbar_n(0,W_i)$ using
treatment-specific prediction datasets. To use it for IPTW, estimate $g_{n,1}(W_i)$
and form inverse probability weights.
\end{estimationbox}
